\documentclass{article}
\usepackage[utf8]{inputenc}
\usepackage[english]{babel}
\usepackage{hyperref}
\usepackage{listings}
\usepackage{xcolor}
\usepackage{amssymb}
\usepackage{graphicx}
\usepackage{pdfpages}
\usepackage{amsopn}
\usepackage{amsmath}

\title{marketpower2: An Interbank model with change of lender mechanism}
\author{Héctor Castillo Andreu}
\date{May 2026}

\begin{document}

\maketitle

\section{Auxiliary files}
\begin{itemize}
    \item \texttt{pyproject.toml}: UV environment configuration and also list of the necessary python packages
\end{itemize}

\section{Interbank model}

\begin{itemize}
    \item \texttt{interbank.py}: use to execute standalone the Interbank simulation.
    \begin{itemize}
        \item It accepts command line options. For instance:
            \begin{lstlisting}[language=bash, basicstyle=\ttfamily\small]
interbank.py --log DEBUG --n 150 --t 2000
interbank.py --save results.gdt p=0.4 param=X
            \end{lstlisting}
        \item When it is used as a package, the sequence should be:
            \begin{lstlisting}[language=Python, basicstyle=\ttfamily\small]
import interbank
model = interbank.Model()
model.config.configure(param=x)
model.forward()
            \end{lstlisting}
    \end{itemize}
    \item Basic options:
        \begin{lstlisting}[language=bash, basicstyle=\ttfamily\small]
        # To list all options:
        interbank.py --help

        # To run a simulation based on exp_runner:
        uv run experimments\exp_min_p_0_1.py --do
        # Same but using only surviving banks to  
        # determine how last the system to die:
        python -m experiments.exp_surviving_4.py --do
        \end{lstlisting}

    \item \texttt{interbank\_lenderchange.py}: It contains the algorithm that control the change of lender in the model.
    \item \texttt{interbank\_statistics.py}: Generates the stats of the model.
    \item \texttt{interbank\_log.py}: Generates the logging.
    \item \texttt{interbank\_testclass.py}: It contains the algorithm that control the change of lender in the model.
    \item \texttt{interbank\_web.py}: Minimal web server for interbank using Flask.
    \item \texttt{exp\_runner.py}: A prototype for executing experiments with different parameters and using MonteCarlo (using concurrent.futures to allow multiple threads).
    \item \texttt{exp\_runner\_surviving.py}: A derivation of the former prototype using ray library to execute in a cluster.
    \item \texttt{experiments/}: directory with all the experiments conducted. The results of that executions are stored in a folder determined inside each experiment.
    \item \texttt{doc/algorithm.pdf}: the PDF schema of the algorithm used in the model to propagate shocks and to balance sheets.
    \item \texttt{experiments/}: directory with all the experiments conducted. The results of that executions are stored in a folder determined inside   each experiment.
    \item \texttt{output/}: directory with the output files (default one, you can change it using the `--output` option).



\end{itemize}


\section{Basic usage of the model}

\begin{figure}[htb]
  \centering
  \includegraphics[width=\textwidth,keepaspectratio]{doc/alg-1}
  \caption{Sequence of steps: grey boxes indicates moments in which that statistic is obtained}
  \label{fig:algorithm}
\end{figure}

\begin{itemize}
    \item \texttt{interbank.py seed=1234 T=500 --p 0.2}: Execute the model with $T=500$ and $LenderChange$ algorithm of $ShockedMarket3$ with an Erdös-Réni with probability of attachment $p_a=0.2$ and using a seed for generating random values of $1234$ (same results if you generate again with other equal parameters and repeat this integer number for seed).
    \item \texttt{interbank.py --save result --output\_format csv}: Save the results in \texttt{result.csv} in $CSV$. With \texttt{--log DEBUG --logfile result.txt} will save also the log in a file.
\end{itemize}


\begin{itemize}
    \item \texttt{interbank.py --normalize\_ir}: Execute the model normalization the interest rate to a range $[0..1]$.
    \item \texttt{interbank.py --robust\_ir}: Execute the model with a similar to a Robust Scaler algorithm normalization.

    In our implementation, robust\_ir is a quantile-based min–max normalization applied per time step, not a true RobustScaler. Here is the formulation:
            \begin{enumerate}
                \item Let determine first the set of interest rates at time \(t\):
                \[
                \mathcal{I}_t = \left\{ r_i \,\middle|\, r_i \in \mathbb{R},\ i = 1,\dots,N_t \right\}
                \]

                \item We compute the low and high quantiles of this set (\(p_{\text{low}} = 5\%\) and \(p_{\text{high}} = 95\%\) by default):
                \[
                q_{\text{low}} = Q_{p_{\text{low}}}(\mathcal{I}_t),
                \qquad
                q_{\text{high}} = Q_{p_{\text{high}}}(\mathcal{I}_t).
                \]

                \item We clip the outliers:

                \[
                \tilde{r}_i = \max\!\left(q_{\text{low}},\, \min\!\left(q_{\text{high}},\, r_i\right)\right)
                \]

                \item And min-max normalize into  \([0, R]\) (with $R=1.0$ by default):

                \[
                r_i^{(\text{robust})} =
                \begin{cases}
                R \, \dfrac{\tilde{r}_i - q_{\text{low}}}{q_{\text{high}} - q_{\text{low}}}, & q_{\text{high}} > q_{\text{low}} \\[6pt]
                0, & q_{\text{high}} = q_{\text{low}}
                \end{cases}
                \]

            \end{enumerate}
           
\end{itemize}

\section{Web server}


A Web interface is available to run the model and visualize the results in charts. It is implemented with Flask and can be launched with \texttt{interbank.py --web} option.
With  \texttt{--web\_mode dashboard --web\_port 8080} we obtain the same as previous command but explicitly setting the mode and port. Available modes are:
\begin{itemize}
    \item \texttt{simulate}: single execution UI.
    \item \texttt{multiple}: parameter multiple-run UI.
    \item \texttt{dashboard}: combined UI with tabs, includes charts for simulate/multiple.
\end{itemize}
In dashboard, chart metric visibility is selectable with checkboxes. Metric selections are persisted in browser `localStorage`.
Backend routes are:
\begin{itemize}
    \item \texttt{POST /api/simulate} for single runs.
    \item \texttt{POST /api/multiple} for multiple parameter runs.
\end{itemize}
The dashboard template file is \texttt{templates/template\_simulation.html}.

\section{Statistics}

Different statistics can be obtained after running the model, either in \textbf{csv} output, or in \textbf{gdt} (Gretl format).
This statistics collect data in each time for the average or individually, depending on the usage. Possible statistics obtained from the model are:

\begin{itemize}
    \item \textbf{asset\_i}: Assets of the lender of this bank ($D + E$)
    \item \textbf{asset\_j}: Assets of the borrowers of this bank ($D + E$)
    \item \textbf{bad\_debt}: Sum of the bad debt
    \item \textbf{bankruptcies}: Number of banks that failed in this step
    \item \textbf{capacity}: Lender capacity ($1 - \frac{E}{E_{\max}}$) of the bank
    \item \textbf{communities}: Subsets of nodes with higher internal edge density than connections to the rest of the graph
    \item \textbf{communities\_not\_alone}:  Number of \textbf{communities} that are not formed by only one node
    \item \textbf{d1,d2}: Sum of demand of loan of borrowers not satisfied by their own capital after shock1 (d1) and shock2 (d2)
    \item \textbf{deposits}: Sum of deposits $D$ of banks (of their balance $L + C + R = D + E$)
    \item \textbf{equity}: Sum of equity $E$ of all banks: $L + C + R = D + E$ (after repayments)
    \item \textbf{fitness}: Fitness ($\mu$) of the bank
    \item \textbf{gcs}: When we use an Erdös–Rényi graph, the Giant Component Size is the largest number of nodes that are interconnected.
    \item \textbf{grade\_avg}: Average number of edges (connections) for the total banks
    \item \textbf{ir}: Interest rate $r$ of the bank who are potential lenders ($\Delta D>0$)
    \item \textbf{ir\_avg}: Average interest rate $r$ of the bank who have really loans with borrowers
    \item \textbf{liquidity}: Total liquidity $L$ of the Banks  $L + C + R = D + E$
    \item \textbf{loans}: Total amount borrowed by the banks
    \item \textbf{num\_loans}:  Num of loans. Value for \textbf{stats\_market} and normal one will be the same
    \item \textbf{num\_of\_rationed}: Number of banks that were rationed in this step (needed money and were without any possible lender)
    \item \textbf{potential\_lenders}: Number of banks in the first shock having a positive shock ($ \Delta D$)
    \item \textbf{prob\_bankruptcy}: Probability of bankruptcy $p_b=1-\frac{E}{E_{max}}$, between $[0..1]$
    \item \textbf{profits}: Profits obtained in that step
    \item \textbf{psi}: Power market ($psi$) of the banks who are potential lenders, value $[0..1]$
    \item \textbf{rationing}: Total amount of the loans $l$ of the banks
    \item \textbf{reserves}: Reserves $R$ in the balance $L + C + R = D + E$
    \item \textbf{var\_D}: Sum of $ \Delta D1$ and $\Delta D2$ 
    \item \textbf{var\_D1}: Amount of ($ \Delta D$) for the bank in first shock
    \item \textbf{var\_D2}: Amount of ($ \Delta D$) for the bank in second shock    

%         \item \textbf{active\_borrowers}: Number of banks that are involved in a loan as borrowers. Both values in global and \textbf{stats\_market} will be the same.
%     \item \textbf{active\_lenders}: Number of banks that are involved in a loan as borrowers. Both values in global and \textbf{stats\_market} will be the same.
% \item \textbf{equity\_lenders}: Average $E$ of banks who are lenders (after loans)
%     \item \textbf{var\_deposits}: Sum of $ \Delta D$ for both shocks
%     \item \textbf{l\_equity}: Log of equity ($\log(E)$)
%     \item \textbf{leverage}: Financial leverage ($l/E$) of the bank considering only the banks that are inside a loan, named \textbf{leverage\_} in Gretl due to name restrictions of the environment.
%     \item \textbf{max\_e}: $E_{\max}$ of the system
%     \item \textbf{max\_lambda}: $\lambda_{\max}$ of the system (only for potential lenders)
%     \item \textbf{num\_banks}: Number of banks currently surviving in the model (interesting when \textbf{allow\_replacement\_of\_bankrupted=False})
%     \item \textbf{policy}: Policy recommendation $\eta$ of the system in the range $[0..1]$. As $\eta$ is a global value, the same number applies for all banks.
%     \item \textbf{potential\_credit\_channels}:  Considering there is a graph of connections between banks, then \textbf{number\_of\_edges()} in the graph
%     \item \textbf{psi\_effective}: Power market ($psi$) of the banks who are really lenders with active loans, value $[0..1]$
%     \item \textbf{real\_t}: Times in which are no loans are removed in the extra statistics generated when we use \textbf{--stats\_market}. Real $t$ instants of time are stored in this variable to track when were really those values are obtained in the original statistics.
%     \item \textbf{systemic\_leverage}: Financial leverage but considering in the mean the total banks of the model $N$

\end{itemize}

The different statistics of information obtained in table~\ref{table1} are classified as:
\begin{itemize}
        \item Global: using \textbf{--save filename}: each data column in \textbf{filename.gdt}  will be obtained for all the $N$ banks in the model for all instants time $T$ (rows)
        \item With \textbf{--stats\_market} what we obtain will be a file as \textbf{filenameb.gdt} for the subsets of banks that in each time are engaged in a real loan. So if in the time $t$ there are no loans, it is removed from this statistics. The special value \texttt{real\_t} indicates which was the original time.
        \item Individual is data that makes sense to be obtained for each bank.
        \item Graphs are data obtained also in \textbf{filename.gdt}, but only we have a \textbf{LenderChange} algorithm with a random graph.
\end{itemize}

\begin{table}[h]
\centering
\begin{tabular}{|l|c|c|c|c|c|}
\hline
Name & Type & Global & \textbf{stats\_market} & Individual & Graphs \\
%\textbf{active\_borrowers} & integer & $\checkmark$ &  $\checkmark$  &  & \\
%\textbf{active\_lenders}   & integer & $\checkmark$ & $\checkmark$  &  & \\
\textbf{asset\_i}          & float & $\overline{x}/0 $ & $\overline{x}/0$ & $\checkmark$ &  \\
\textbf{asset\_j}          & float & $\overline{x}/0 $ & $\overline{x}/0$ & $\checkmark$ &  \\
\textbf{bad\_debt}         & float & $ \sum$ &  $ \sum$ & $\checkmark$ & \\
\textbf{bankruptcies}      & integer & $ \sum$ &  $ \sum$ & $\checkmark$ &  \\
%\textbf{bankrupcty\_rationed} & integer & $ \sum$ &  $ \sum$ &  &  \\
%\textbf{best\_lender}      &  integer &  $\checkmark$ & $\checkmark$  &  & \\
%\textbf{best\_lender\_clients}  & integer &  $\checkmark$ &  $\checkmark$ &  & \\
\textbf{capacity}                 & float &  $\overline{x}/nan $ & $\overline{x}/nan $ &  &  \\
\textbf{communities} & integer &  &   &  & $\checkmark$\\
\textbf{communities\_not\_alone} & integer &  &   &  & $\checkmark$ \\
\textbf{d1,d2}            & float &  $\sum$ & $\sum$ &  &  \\
\textbf{deposits}          & float & $ \sum$ &  $ \sum$ & $\checkmark$ &  \\
\textbf{equity}            & float &  $ \sum$ &  $ \sum$ &  & \\
%\textbf{equity\_lenders}            & float &  $\overline{x}$ &  &  & \\
\textbf{fitness}           & float & $\overline{x}$ & $\overline{x}/nan $ &  $\checkmark$  & \\
\textbf{gcs}               & integer &  &   &  & $\checkmark$ \\
\textbf{grade\_avg}        & integer &  &   &  & $\checkmark$\\

\textbf{ir,ir\_avg}                &  float &  $\overline{x}/0$ & $\overline{x} / nan $  & $\checkmark$ &  \\
%\textbf{l\_equity}         & float &  $ \sum$ &  $ \sum $ & & \\
%\textbf{leverage} / \textbf{leverage\_}         & float  &  $\overline{x}$ & $\overline{x}/nan $  & $\checkmark$  &  \\
\textbf{liquidity}         &  float &  $\sum$ &  $\sum$ &  $\checkmark$ &    \\
%\textbf{max\_e}           &  float &  $\checkmark$ & $\checkmark$ &  & \\
%\textbf{max\_lambda}      &  float &  $\checkmark$ & $\checkmark$ &  & \\
%\textbf{num\_banks}        &  integer &  $\checkmark$ & $\checkmark$ &  & \\
\textbf{loans}        & float &  $ \sum$ &  $ \sum$ & $\checkmark$ &  \\
\textbf{num\_loans}        &  integer &  $ \checkmark$ &  $\checkmark$ & $\checkmark$ &  \\
\textbf{num\_of\_rationed} & integer &  $\checkmark$ & $\checkmark$ &  & \\
%\textbf{policy}            & float &  $\checkmark$ & $\checkmark$ &  & \\
%\textbf{potential\_credit\_channels} & integer &  $\checkmark$ & $\checkmark$ &  & \\
%\textbf{potential\_lenders} & integer &  $\checkmark$ & $\checkmark$ &  & \\
\textbf{prob\_bankruptcy}  &  float &  $\overline{x}/nan$ & $\overline{x}/nan$  &  $\checkmark$  & \\
\textbf{profits}           & float &  $ \sum$ &  $ \sum$ &  & $\checkmark$ \\
\textbf{psi} &  float &  $\overline{x}/0 $ & $\overline{x}/nan $ &  & $\checkmark$ \\
\textbf{rationing}         &  float &  $\sum$ &  $\sum$ &  & $\checkmark$  \\
%\textbf{real\_t}         &  integer &  &  $\checkmark$ &  &   \\
\textbf{reserves}          & float & $ \sum$ &  $ \sum$ &  & $\checkmark$ \\
\textbf{var\_D,var\_D1,var\_D2} &  float &   $\sum$  &  $\sum$  &  &  $\checkmark$  \\
%\textbf{systemic\_leverage} & float &  $\overline{x}$ & $ \overline{x} $  &  &  \\
\hline
\end{tabular}
\caption{ Legend for the different columns are: $\checkmark$=value without any modification. $\sum$=sum of the value for all banks. $ \overline{x}$ = average of the value for all banks. $0$ = No banks in this statistic. $nan$=Instead of zero, the value of "not a number" is used}
\label{table1}
\end{table}
\end{document}